\chapter*{Resumen}
\addcontentsline{toc}{chapter}{Resumen}
\markboth{Resumen}{Resumen}

La enseñanza del Método de los Elementos Finitos (MEF) enfrenta una brecha pedagógica persistente: los textos clásicos exponen la formulación matricial con rigor pero en abstracto, y el software profesional opera como una caja negra que oculta el procedimiento que el estudiante necesita comprender. Esta tesis aborda esa brecha como investigación basada en diseño y reporta su primer ciclo: el desarrollo de EduFEM, una herramienta educativa interactiva en español para el análisis estructural de problemas bidimensionales de elasticidad lineal en tensión plana y deformación plana, entendido como el análisis tenso-deformacional de medios continuos, y la evaluación de su potencial didáctico. La hipótesis sostiene que exponer de forma transparente, interactiva y verificable el canal de cálculo del MEF sobre el modelo del propio estudiante apoya la comprensión de los fundamentos del método. La herramienta combina un motor de cálculo con elementos isoparamétricos cuadriláteros Q4 y Q9, una interfaz gráfica que integra pre-proceso, proceso y post-proceso sobre un único lienzo de malla, ocho módulos educativos que exponen paso a paso cada etapa del cálculo, del mapeo isoparamétrico al ensamblaje global, y una memoria de cálculo automática con sustitución numérica que el estudiante puede reproducir a mano. La corrección del motor se sometió a una batería de verificación y validación: con el método de soluciones manufacturadas, en ambos estados planos y sobre mallas regulares y distorsionadas, se confirmaron las tasas de convergencia asintóticas teóricas del desplazamiento —orden dos en norma $L^2$ y orden uno en seminorma $H^1$ para Q4; tres y dos para Q9—; la viga de Timoshenko arrojó errores del 0,04\,\% en la tensión normal y del 0,26\,\% en la flecha frente a la solución analítica, y diferencias de hasta el 0,56\,\% frente a SAP2000; y la membrana de Cook evidenció el bloqueo por cortante (\emph{shear-locking}) del Q4 frente a la convergencia rápida del Q9, que a igual número de grados de libertad (GDL) reduce el error en un orden de magnitud. El potencial didáctico se evaluó con tres instrumentos documentales fijados a priori: una tabla de especificaciones, que verificó la cobertura de las once unidades de contenido del canal de cálculo en los niveles cognitivos de comprender y aplicar; un mapa de conjeturas con su matriz de trazabilidad, que verificó que cada uno de los diez principios de aprendizaje que fundamentan el diseño tiene una encarnación comprobable en el software; y una evaluación heurística con el instrumento LORI, que documentó la evidencia disponible en sus nueve dimensiones. La efectividad —la comprensión lograda por los estudiantes— queda como conjetura teórica del mapa, con la prueba de campo del segundo ciclo diseñada en las recomendaciones.

\noindent\textbf{Palabras clave:} método de elementos finitos, software educativo, investigación basada en diseño, elementos isoparamétricos Q4/Q9, verificación y validación, memoria de cálculo
